\section{局限与结论}\label{sec:limits}

\subsection{局限}\label{sec:limitations}

\textbf{（1）单一检查点、单一规模，且只训练到 20\%。}本文全部数字来自
\texttt{gridman\_medium\_v\_0d3\_pretrain}（437.8M 参数，step $144{,}000/705{,}600$），
没有验证集、没有同预算基线（\S\ref{sec:b3} 未完成）。
因此：绝对值（如 $1.08\bits$）\textbf{不可}与任何外部模型对照，
所有结论只在\textbf{条件间}成立；「有效记忆 $1$--$4$k token」也\textbf{不能}排除
「换个训练更充分的检查点会显著更长」的可能。加验证集、训完 medium、补一个同预算的
byte Transformer 基线，是本稿升级为投稿稿的最低条件。

\textbf{（2）精确检索与离散绑定。}合成键值对与词对联想上，本检查点的绑定能力弱到接近机会线
（\S\ref{sec:gline}）；据此本文把 needle 类精确检索降级为限定 claim，
并如实说明「统计/语义级的长程记忆」才是本文证据支持的范围。B4 预实验本身\textbf{尚未运行}。

\textbf{（3）热启动与训练语料的污染边界。}状态级样本记忆 $=-0.714\bits$
说明热启动起点含着训练语料的具体痕迹（\S\ref{sec:dline}）。
本文据此要求主表双报起点状态；但「双报」只解决报告透明，
不解决\textbf{污染本身}——严格版本需要用完全无重叠的评测集重做 A/B/C 三组。
$\mathrm{n}$-gram 重叠率的精确统计 \pending{待补，见附录 \ref{app:decontam}}。

\textbf{（4）受控遗忘曲线缺失。}计划中的 C 线（注入键值对、按距离与组件分解测 recall 衰减）
因实现缺陷作废（附录 \ref{app:c1}），修订版尚未重跑。
这意味着本文\textbf{没有}一条真正受控的遗忘曲线：
\S\ref{sec:b2} 的重复距离是对自然文本的观察性测量，
其各距离桶之间并非完全受控（不同距离上重复的是不同 $\mathrm{n}$-gram 群体，首次出现难度本就不同）。

\textbf{（5）实现版本已演进，部分归因结论绑定旧语义。}本文数字全部出自 v\_0d3 检查点；
现行 v\_0d4 把 \texttt{memory} 改为序列级私有槽位（无跨流平均、无跨卡同步）。
推理侧（$\mathrm{batch}=1$）两版逐位对应，主结论（S1/S2、A1/A2/B2）不受影响；
但\textbf{训练侧语义已变}：v\_0d3 的「重置 memory $=$ 动共享知识」在 v\_0d4 中退化为只影响单流，
\S\ref{sec:s1-attr} 的组件归因、A5 的 memory 回滚语义在 v\_0d4 检查点上需重新验证。
此外 v\_0d3/v\_0d4 存档互不兼容，跨版续训不可行。

\textbf{（6）未完成项。}A4（状态错配）、B3（等效滑窗基线）、B4（needle 预实验）、
F 线（训练侧全部实验）、E3（精度链漂移）、E4（自由运行分叉）均未运行，
正文与附录中一律标注为待补，不以任何形式给出推测数值。

\textbf{（7）口径订正史（主动披露）。}本稿在定稿前对实验包做了逐脚本审读，
修正了三处口径问题（$\mathrm{n}$-gram 分桶缺陷导致的一条错误结论已撤回、
A1 计分窗口未对齐、C1 答案位置偏移导致该实验作废），并补了模型无关的回归自检。
保留这段历史的原因是：其中一条（「长程零衰减」）曾经看起来非常漂亮，
而它恰恰来自一个会让所有桶报告同一份数据的实现缺陷。\textbf{分桶/配对类统计必须断言样本数互不相同}，
这条纪律比结论本身更重要。

\subsection{结论}\label{sec:conclusion}

本文把「固定大小的持续状态能承载多少跨上下文信息」当作一个可测量的经验问题，
在一个 437.8M 参数的 byte 级有状态语言模型上，用 teacher-forced bits/byte 与状态手术
（热启动/空白、commit/$\mathit{lock}$、分量重置/冻结）做了系统消融，得到三条可以直接引用的结论。

\textbf{第一，状态确实在携带信息，且写入有因果贡献。}同一计分窗口下热启动比空白低 $3.73\bits$
且无收敛趋势；锁定跨 patch 写入的代价是 $+1.00\pm0.09\bits$，并随上下文累积。S1 成立。

\textbf{第二，承载者是块级循环状态，不是参数量最大的矩阵记忆。}冻结 \texttt{recurrent}
的代价（$+1.022\bits$）与冻结全部写入（$+1.034\bits$）几乎相同；
冻结 delta 规则矩阵记忆的代价为零（$-0.0005\bits$），
但去掉热启动使损失升高 $3.7\bits$ 以上——\textbf{矩阵记忆的价值在读取训练稳态先验，
不在增量适应}。这个结论由 A2/A5/A6/B2-freeze 四组独立实验交叉支撑。

\textbf{第三，有效记忆距离约为 $1$--$4$k token，且长度不是瓶颈。}自然重复的复现增益集中在
$<\!1$k（$+0.084\nats$），$4$k 之外衰减到零；同时 $512$\,KB 连续写入下状态范数有界
（memory $\times1.13$），重置频率扫描单调无 U 形。因此「永续写入」在数值上是安全的，
瓶颈在容量与写入/读出机制，而不是「能不能一直写下去」。

\textbf{对「永续上下文」主张的净判断}：固定大小状态\textbf{可以}持续携带信息，
并且\textbf{可以}在远长于其有效记忆距离的流上稳定运行；
但在本文这个检查点上，它的记忆能力等效于数 k token 的滑动记忆，
尚不足以支撑「记住了整条流」的强主张。
下一步的判决性实验是 BPTT 窗口扫描（$\mathrm{bptt}=1$ 保留多少长程能力）
与永续 vs 边界重置训练的对照（\S\ref{sec:f-line}）——
前者区分「前向通道上界」与「梯度通道贡献」，后者直接回答「训练时不清状态是收益还是负担」。
